\section{导读：一场发生在收购前的“最后访谈”}

本节先说明为什么这期值得单独整理。它录制于 2025 年 12 月 1 日。录制时，Meta 尚未宣布全资收购 Manus；节目发布时，它已经成了 Manus 的最后访谈。这个时间差很重要，因为它让访谈呈现出一种罕见状态：嘉宾在不知道结局的情况下，完整讲出了自己的创业路径、产品判断、增长逻辑、行业观点和内心隐忧。

本期主线不是“某公司被收购”的新闻复盘，而是一个 Agent 产品团队如何从个人开发、NLP、知识图谱、Open IE、GUI Agent、Manus 增长，再走到“AI 更像制造业”和“害怕复杂”的产品哲学。它和 EP130 的 Dokie、EP139 的 Agent 综述、EP129 的模型公司上市视角可以连成一条线：模型能力正在进入产品和公司组织层面，真正难的是把能力变成可持续使用、可持续收入和可持续简单的系统。

\begin{importantbox}{本期核心命题}
Manus 的故事不是单纯“Agent 爆款”，而是把 AI 能力包装成用户愿意持续使用和付费的生产力系统。它的最大风险不只是大厂竞争，而是增长过程中失去特色、变得复杂。
\end{importantbox}

\begin{knowledgebox}{视觉策略说明}
本视频是固定访谈画面，没有 slides、白板或产品演示。正文不重复插入人物帧；封面用于来源识别，正文用 Agent 漂流路线、制造业隐喻、复杂性风险和市场地图等概念图解释内容。
\end{knowledgebox}

\subsection{本章小结}

EP128 是一个“事后带着结局读，会更有意义”的访谈。它不仅讲 Manus 怎么增长，也讲一个 Agent 公司怎样理解平台、产品、组织、复杂性和未来竞争。

